% ============================================================
% Klausur-Formelsammlung: BWL / VWL
% DHBW -- Technik & Produktion, Semester 2
% ============================================================
\documentclass[10pt,a4paper]{article}
\usepackage[a4paper, top=1.5cm, bottom=1.5cm, left=1.5cm, right=1.5cm]{geometry}
\usepackage[ngerman]{babel}
\usepackage[utf8]{inputenc}
\usepackage[T1]{fontenc}
\usepackage{amsmath}
\usepackage{amssymb}
\usepackage{booktabs}
\usepackage{tabularx}
\usepackage{array}
\usepackage{multirow}
\usepackage{multicol}
\usepackage{xcolor}
\usepackage{mdframed}
\usepackage{enumitem}
\usepackage{fancyhdr}
\usepackage{titlesec}
\usepackage{lmodern}
\usepackage{microtype}
\usepackage{colortbl}
\usepackage{tcolorbox}
\tcbuselibrary{skins,breakable}

% ---- Farben ----
\definecolor{bwlblue}{RGB}{0,60,120}
\definecolor{vwlgreen}{RGB}{0,100,50}
\definecolor{formulabg}{RGB}{232,242,255}
\definecolor{hinweisbg}{RGB}{255,248,220}
\definecolor{vwlbg}{RGB}{230,248,238}
\definecolor{rowbg}{RGB}{220,234,248}
\definecolor{rowvwl}{RGB}{210,240,220}

% ---- Section-Format ----
\titleformat{\section}{\large\bfseries\color{bwlblue}}{}{0em}{}[\color{bwlblue}\titlerule]
\titleformat{\subsection}{\normalsize\bfseries\color{bwlblue!80!black}}{}{0em}{}
\titleformat{\subsubsection}{\small\bfseries\itshape}{}{0em}{}
\titlespacing*{\section}{0pt}{5pt}{3pt}
\titlespacing*{\subsection}{0pt}{4pt}{2pt}
\titlespacing*{\subsubsection}{0pt}{2pt}{1pt}

% ---- tcolorbox-Stile ----
\newtcolorbox{formelbox}{colback=formulabg, colframe=bwlblue, boxrule=0.6pt,
  left=4pt, right=4pt, top=3pt, bottom=3pt}
\newtcolorbox{hinweisbox}{colback=hinweisbg, colframe=orange!70!black, boxrule=0.5pt,
  left=4pt, right=4pt, top=2pt, bottom=2pt, fonttitle=\small\bfseries}
\newtcolorbox{vwlbox}{colback=vwlbg, colframe=vwlgreen, boxrule=0.6pt,
  left=4pt, right=4pt, top=3pt, bottom=3pt}

% ---- Tabellenspalten ----
\newcolumntype{L}[1]{>{\raggedright\arraybackslash}p{#1}}
\newcolumntype{C}[1]{>{\centering\arraybackslash}p{#1}}

% ---- Seitenkopf ----
\pagestyle{fancy}
\fancyhf{}
\fancyhead[L]{\small\textbf{Formelsammlung BWL/VWL} -- DHBW}
\fancyhead[R]{\small Seite \thepage}
\renewcommand{\headrulewidth}{0.4pt}
\setlength{\parindent}{0pt}
\setlength{\parskip}{2pt}
\setlength{\columnsep}{8pt}
\sloppy
\hfuzz=4pt
\hyphenation{Wie-der-be-schaf-fungs-kos-ten An-schaf-fungs-kos-ten
  For-de-rungs-aus-fall-wag-nis Ge-währ-leis-tungs-wag-nis
  Kos-ten-stel-len-rech-nung Kos-ten-trä-ger-rech-nung
  Kos-ten-ar-ten-rech-nung Be-triebs-ab-rech-nungs-bo-gen
  Haupt-re-fi-nan-zie-rungs-ge-schäft Preis-elas-ti-zi-tät
  Pro-duk-ti-ons-fak-to-ren Ge-winn-ver-tei-lung
  Ei-gen-ka-pi-tal Fremd-ka-pi-tal Ge-samt-ka-pi-tal
  Ein-zel-un-ter-neh-men Ge-sell-schaft-er
  Min-dest-re-ser-ve-satz Geld-schöp-fungs-mul-ti-pli-ka-tor
  Mehr-li-ni-en-or-ga-ni-sa-ti-on Ein-li-ni-en-or-ga-ni-sa-ti-on}

% ============================================================
\begin{document}

\begin{center}
  {\LARGE\bfseries\color{bwlblue} Klausur-Formelsammlung: BWL / VWL}\\[4pt]
  {\normalsize DHBW Technik \& Produktion \quad|\quad Kosten- und Leistungsrechnung + Volkswirtschaftslehre}\\[2pt]
  {\small\color{gray} Klausur: 4--5 Aufgaben, 100 Punkte, 50\% BWL / 50\% VWL \quad|\quad Schwerpunkt BWL: 3 Stufen der KLR}
\end{center}
{\color{bwlblue}\hrule height 1pt}
\vspace{4pt}

% ============================================================
%  SEITE 1 -- ÜBERBLICK KLR & ABGRENZUNGSRECHNUNG
% ============================================================
\section{1 \quad Aufbau der KLR und Abgrenzungsrechnung}

\begin{multicols}{2}

\subsection{Die 3 Stufen der Kostenrechnung}

{\small
\begin{tabular}{|L{2.3cm}|L{1.7cm}|L{3.2cm}|}
\hline
\rowcolor{rowbg}
\textbf{Stufe} & \textbf{Frage} & \textbf{Instrument / Aufgabe} \\
\hline
1. Kostenartenrechnung & Welche Kosten? & Kostenarten\-liste \\
\hline
2. Kostenstellenrechnung & Wo angefallen? & BAB (Betriebsabrechnungsbogen) \\
\hline
3. Kostenträgerrechnung & Wofür? & Kalkulation \\
\hline
\end{tabular}
}

\vspace{4pt}
\textbf{Wichtige Kostenbegriffe:}

\begin{tabular}{|L{2.5cm}|L{4.5cm}|}
\hline
\rowcolor{rowbg}
\textbf{Begriff} & \textbf{Bedeutung} \\
\hline
Einzelkosten & Direkt einem Produkt (Kostenträger) zurechenbar, z.\,B. Rohstoff, Fertigungslohn \\
\hline
Gemeinkosten & Nicht direkt zurechenbar; Verteilung über Zuschlagssätze \\
\hline
Fixkosten $K_f$ & Unabhängig von der Ausbringungsmenge $x$ \\
\hline
Variable Kosten $K_v$ & Proportional zur Menge $x$ \\
\hline
Kalk. Kosten & Nur in der KLR (kein Aufwand in FIBU) \\
\hline
Neutrale Aufwendungen & Nur in FIBU, nicht in KLR \\
\hline
\end{tabular}

\columnbreak

\subsection{Abgrenzungsrechnung: FIBU $\to$ KLR}

\textbf{Überleitung Aufwand zu Kosten:}

\begin{formelbox}
\begin{align*}
\text{Kosten} &= \text{Grundkosten} - \text{Neutraler Aufwand}\\
              &\quad + \text{Kalk. Kosten}
\end{align*}
{\footnotesize
\textbf{Grundkosten} = in FIBU \textit{und} KLR \quad
\textbf{Neutraler Aufwand} = nur FIBU\\
\textbf{Kalk. Kosten} = nur KLR}
\end{formelbox}

\vspace{4pt}
\begin{tabular}{|L{4cm}|L{3cm}|}
\hline
\rowcolor{rowbg}
\textbf{Nur FIBU $\to$ herausrechnen} & \textbf{Nur KLR $\to$ hinzurechnen} \\
\hline
Betriebsfremder Aufwand & Kalk. Unternehmerlohn \\
(Spenden, Verluste aus Beteiligungen) & \\
\hline
Außerordentl. Aufwand (Brand, Unfall) & Kalk. Miete (Eigengebäude) \\
\hline
Periodenfremder Aufwand & Kalk. Zinsen (auf EK) \\
\hline
Buchm. AfA (falls abweichend) & Kalk. Abschreibungen \\
\hline
-- & Kalk. Wagnisse \\
\hline
\end{tabular}

\vspace{4pt}
\begin{hinweisbox}
\textbf{Prüfungshinweis:} Buchm. Abschreibungen in FIBU $\ne$ kalk. Abschreibungen in KLR. Wenn beide existieren: buchm. AfA heraus, kalk. AfA hinein. Neutrale Erträge analog behandeln (z.\,B. Mieteinnahmen aus betriebsfremdem Gebäude).
\end{hinweisbox}

\end{multicols}

\newpage

% ============================================================
%  SEITE 2 -- KALKULATORISCHE KOSTEN
% ============================================================
\section{2 \quad Kalkulatorische Kosten}

\begin{multicols}{2}

\subsection{Kalk. Abschreibungen}

Ziel: tatsächlichen Werteverzehr der Anlagen erfassen -- unabhängig von Steuerrecht.

\subsubsection{Lineare Abschreibung}
\begin{formelbox}
\[
\text{Kalk. AfA} = \frac{\text{Bewertungsansatz}}{\text{Nutzungsdauer (ND)}}
\]
\end{formelbox}

\subsubsection{Geometrisch-degressive Abschreibung}
\begin{formelbox}
\[
\text{AfA}_t = \text{RBW}_{t-1} \times p\%
\]
\[
\text{RBW}_t = \text{RBW}_{t-1} - \text{AfA}_t
\]
\end{formelbox}

\begin{tabular}{|L{2cm}|L{5cm}|}
\hline
RBW & Restbuchwert zu Beginn des Jahres $t$ \\
$p\%$ & Abschreibungsprozentsatz \\
\hline
\end{tabular}

\subsubsection{Bewertungsgrundlage}
\begin{tabular}{|L{3.2cm}|L{3.6cm}|}
\hline
\rowcolor{rowbg}
\textbf{Basis} & \textbf{Verwendung} \\
\hline
AK (Anschaffungs\-kosten) & Konservativ; nach FIBU-Recht \\
WBK (Wieder\-beschaffungs\-kosten) & Substanzerhaltung; in KLR bevorzugt \\
\hline
\end{tabular}

\vspace{4pt}
\begin{hinweisbox}
Bei Anwendung der linearen AfA: betriebsnotwendiges AV oft als \textbf{halbe AK} angesetzt (= Durchschnittswert): $\text{AV}_\varnothing = \text{AK}/2$
\end{hinweisbox}

\subsection{Kalk. Zinsen}

Opportunitätskosten des im Betrieb gebundenen Kapitals (auch für Eigenkapital!).

\begin{formelbox}
\begin{align*}
\text{BNK} &= \text{AV} + \text{UV} - \text{Abzugskapital}\\
\text{Kalk. Zinsen} &= \text{BNK} \times i
\end{align*}
\footnotesize AV = Anlagevermögen; UV = Umlaufvermögen
\end{formelbox}

\begin{tabular}{|L{2.5cm}|L{4.5cm}|}
\hline
\rowcolor{rowbg}
\textbf{Symbol} & \textbf{Bedeutung} \\
\hline
BNK & Betriebsnotwendiges Kapital \\
AV & Anlagevermögen (ggf. als Durchschnittswert) \\
UV & Umlaufvermögen (Vorräte, Forderungen, Kasse) \\
Abzugskapital & Zinslos überlassenes Fremdkapital: Lieferantenverbindlichkeiten, Kundenanzahlungen, unverzinsliche Rückstellungen \\
$i$ & Kalkulationszinssatz (z.\,B. 8\% p.\,a.) \\
\hline
\end{tabular}

\columnbreak

\subsection{Kalk. Unternehmerlohn}

Vergütung für die persönliche Arbeitsleistung des Unternehmers (Einzelunternehmer, Personengesellschafter erhalten kein Gehalt $\to$ Opportunitätskosten).

\begin{formelbox}
\[
\text{Kalk. Unternehmerlohn} = 18 \times \sqrt{\text{Jahresumsatz in EUR}}
\]
\end{formelbox}

\begin{hinweisbox}
Wenn der Unternehmer z.\,B. 150.000\,€ Jahresumsatz hat:\\
$18 \times \sqrt{150.000} = 18 \times 387{,}3 \approx 6.971$\,€
\end{hinweisbox}

\subsection{Kalk. Wagnisse}

Risikokosten für typische betriebliche Risiken (nicht versicherte Schäden, Forderungsausfälle, Lagerrisiken).

\begin{formelbox}
\[
\text{Kalk. Wagnis} = \text{Zuschlagsatz} \times \text{Bezugsgröße}
\]
\end{formelbox}

\textbf{Kalk. Wagnisarten:}

\begin{tabular}{|L{2.5cm}|L{4.5cm}|}
\hline
\rowcolor{rowbg}
\textbf{Wagnis} & \textbf{Bezugsgröße} \\
\hline
Bestandswagnis & Materialwert (Lager) \\
Forderungsausfall- wagnis & Umsatz / Forderungen \\
Gewährleistungs- wagnis & Herstellkosten der Periode \\
Anlagenwagnis & Anlagevermögen \\
\hline
\end{tabular}

\vspace{4pt}
\begin{hinweisbox}
\textbf{Gesamtübersicht kalk. Kosten in der Klausur:}\\
Gegeben: Buchwert der Maschine, Nutzungsdauer, Zinssatz\\
$\Rightarrow$ Berechne: kalk. AfA (WBK $\div$ ND), dann kalk. Zinsen (BNK $\times i$)\\
Ergebnis fließt in BAB als Gemeinkostenposition ein.
\end{hinweisbox}

\end{multicols}

\newpage

% ============================================================
%  SEITE 3 -- MATERIALVERBRAUCH & BEWERTUNG
% ============================================================
\section{3 \quad Materialverbrauch: Erfassung und Bewertung}

\begin{multicols}{2}

\subsection{Erfassung des Materialverbrauchs}

\begin{tabular}{|L{2.2cm}|L{1.3cm}|L{1.3cm}|L{2cm}|}
\hline
\rowcolor{rowbg}
\textbf{Methode} & \textbf{Aufw.} & \textbf{Genau.} & \textbf{Prinzip} \\
\hline
Retrograd & gering & gering & $x \times$ Norm \\
Inventur & mittel & mittel & AB $+$ Z. $-$ EB \\
Skontration & hoch & hoch & AB $+$ Z. $-$ Abg. \\
\hline
\end{tabular}

\vspace{4pt}
\textbf{Retrograde Methode:}
\begin{formelbox}
\[\text{Verbrauch} = x \times \text{Sollverbrauch/Einheit}\]
\footnotesize $x$ = hergestellte Stückzahl
\end{formelbox}

\textbf{Inventurmethode} (Endbestand durch körperliche Aufnahme):
\begin{formelbox}
\[\text{Verbrauch (Abgang)} = \text{AB} + \text{Zugang} - \underbrace{\text{EB}}_{\text{gezählt}}\]
\end{formelbox}

\textbf{Skontrationsmethode} (Fortschreibung, genaueste Methode):
\begin{formelbox}
\[\underbrace{\text{EB}}_{\text{errechnet}} = \text{AB} + \text{Zugang} - \text{Abgang}\]
\end{formelbox}
Jede Materialentnahme wird per Entnahmeschein erfasst.

\textbf{Materialarten:}\par
\begin{tabular}{|L{1.9cm}|L{4.5cm}|}
\hline
Rohstoffe & Hauptbestandteil des Produkts (Einzelkosten) \\
Hilfsstoffe & Nebenbestandteil im Produkt (oft Gemeinkosten) \\
Betriebsstoffe & Verbraucht, gehen nicht ins Produkt (Gemeinkosten) \\
\hline
\end{tabular}

\columnbreak

\subsection{Bewertung des Materialverbrauchs}

\textbf{Zu welchem Preis wird verbrauchtes Material bewertet?}

\begin{tabular}{|L{2.2cm}|L{4.8cm}|}
\hline
\rowcolor{rowbg}
\textbf{Methode} & \textbf{Prinzip und Wirkung} \\
\hline
\textbf{FIFO} & \textit{First In First Out}: Älteste Bestände zuerst verbraucht. Endbestand $=$ aktuelle Preise. Bei steigenden Preisen: niedrige Kosten $\to$ höherer Gewinnausweis. \\
\hline
\textbf{LIFO} & \textit{Last In First Out}: Neueste Bestände zuerst. Verbrauch $=$ aktuelle Preise. Bei steigenden Preisen: hohe Kosten $\to$ niedrigerer Gewinnausweis. \\
\hline
\textbf{HIFO} & \textit{Highest In First Out}: Teuerste Bestände zuerst. Endbestand zu \textbf{niedrigsten} Preisen. Höchste Kostenverrechnung. \\
\hline
\textbf{LOFO} & \textit{Lowest In First Out}: Billigste Bestände zuerst. Endbestand zu \textbf{höchsten} Preisen. Niedrigste Kostenverrechnung. \\
\hline
\textbf{Durchschnitt} & Gewogener Ø-Preis aller Zugänge. Periodisch oder permanent. \\
\hline
\end{tabular}

\vspace{4pt}
\textbf{Durchschnittspreis-Formel:}
\begin{formelbox}
\[
\bar{p} = \frac{\sum_{i}(\text{Menge}_i \times \text{Preis}_i)}{\sum_{i} \text{Menge}_i}
\]
\end{formelbox}

\begin{hinweisbox}
\textbf{Rechenbeispiel:}\\
AB: 100\,Stk. à 2\,€; Zugang: 50\,Stk. à 3\,€; Entnahme: 80\,Stk.\\
\textbf{FIFO:} $80 \times 2 = 160$\,€ \quad \textbf{LIFO:} $50\times 3 + 30\times 2 = 210$\,€\\
\textbf{Durchschnitt:} $\bar{p} = \frac{200+150}{150} = 2{,}33$\,€ $\to 80\times 2{,}33 = 186{,}67$\,€
\end{hinweisbox}

\end{multicols}

\newpage

% ============================================================
%  SEITE 4 -- BETRIEBSABRECHNUNGSBOGEN (BAB)
% ============================================================
\section{4 \quad Betriebsabrechnungsbogen (BAB) -- Kostenstellenrechnung}

\begin{multicols}{2}

\subsection{Zweck und Aufbau des BAB}

Der BAB verteilt die Gemeinkosten aus der Kostenartenrechnung verursachungsgerecht auf die Kostenstellen und ermittelt die \textbf{Zuschlagssätze} für die Kalkulation (Stufe~3).

\textbf{Ablauf:}
\begin{enumerate}[noitemsep, leftmargin=*]
  \item Primärkosten (aus FIBU/Kostenartenrechnung) eintragen
  \item Sekundärkosten: Hilfskostenstellen umlegen
  \item Zuschlagssätze berechnen
\end{enumerate}

\vspace{4pt}
\textbf{Typische Kostenstellen im BAB:}

\begin{tabular}{|L{2cm}|L{2.5cm}|L{3cm}|}
\hline
\rowcolor{rowbg}
\textbf{Stelle} & \textbf{Art} & \textbf{Zuschlag auf \ldots} \\
\hline
Material & Hauptkostenstelle & MEK \\
Fertigung & Hauptkostenstelle & Fertigungslöhne (FL) \\
Verwaltung & Hauptkostenstelle & Herstellkosten (HK) \\
Vertrieb & Hauptkostenstelle & Herstellkosten (HK) \\
Gebäude, Kantine, Fuhrpark & Hilfskostenstelle & wird umgelegt \\
\hline
\end{tabular}

\subsection{Zuschlagssätze}

\begin{formelbox}
\[
\text{MGK\%} = \frac{\text{Materialgemeinkosten}}{\text{MEK}} \times 100
\]
\[
\text{FGK\%} = \frac{\text{Fertigungsgemeinkosten}}{\text{FL}} \times 100
\]
\[
\text{VwGK\%} = \frac{\text{Verwaltungsgemeinkosten}}{\text{HK}} \times 100
\]
\[
\text{VtGK\%} = \frac{\text{Vertriebsgemeinkosten}}{\text{HK}} \times 100
\]
\end{formelbox}

\begin{tabular}{|L{2cm}|L{5cm}|}
\hline
\rowcolor{rowbg}
\textbf{Abk.} & \textbf{Bedeutung} \\
\hline
MEK & Materialeinzelkosten \\
FL & Fertigungslöhne (= Fertigungseinzelkosten) \\
HK & Herstellkosten ($= $ MK $+$ FK) \\
MGK & Materialgemeinkosten \\
FGK & Fertigungsgemeinkosten \\
VwGK & Verwaltungsgemeinkosten \\
VtGK & Vertriebsgemeinkosten \\
\hline
\end{tabular}

\columnbreak

\subsection{Umlageverfahren für Hilfskostenstellen}

\textbf{Anbauverfahren} (einfachstes Verfahren):
\begin{itemize}[noitemsep, leftmargin=*]
  \item Hilfskostenstellen direkt auf Hauptkostenstellen umgelegt
  \item Gegenseitige Leistungen zwischen Hilfskostenstellen werden \textbf{ignoriert}
  \item Ungenauer, aber einfach
\end{itemize}

\vspace{4pt}
\textbf{Stufenleiterverfahren} (Treppenumlage):
\begin{itemize}[noitemsep, leftmargin=*]
  \item Hilfskostenstellen werden nacheinander abgerechnet
  \item Reihenfolge: Hilfskostenstelle, die am meisten an andere liefert, zuerst
  \item Bereits abgerechnete Hilfskostenstellen erhalten \textbf{keine} Umlage mehr
  \item Einseitige Berücksichtigung der Leistungsbeziehungen
\end{itemize}

\vspace{4pt}
\textbf{Gleichungsverfahren} (genauestes Verfahren):
\begin{itemize}[noitemsep, leftmargin=*]
  \item Alle gegenseitigen Leistungsbeziehungen vollständig berücksichtigt
  \item Gleichungssystem aufstellen und lösen
\end{itemize}

\begin{formelbox}
\textbf{Beispiel mit 2 Hilfskostenstellen A und B:}\\[4pt]
$K_A = \text{PK}_A + \alpha_{BA} \cdot K_B$\\[2pt]
$K_B = \text{PK}_B + \alpha_{AB} \cdot K_A$\\[4pt]
Gleichungssystem lösen $\Rightarrow$ $K_A$ und $K_B$ einsetzen\\[2pt]
$\alpha_{BA}$ = Anteil, den B an A liefert
\end{formelbox}

\begin{hinweisbox}
\textbf{Klausurhinweis:} Das \textbf{Iterationsverfahren} ist laut Dozent \textbf{nie} in der Klausur vorgekommen. Stufenleiter- und Gleichungsverfahren kennen!
\end{hinweisbox}

\vspace{4pt}
\textbf{Normalzuschlagssätze vs. Istzuschlagssätze:}

\begin{tabular}{|L{2.5cm}|L{4.5cm}|}
\hline
\rowcolor{rowbg}
\textbf{Art} & \textbf{Merkmal} \\
\hline
Ist-Zuschlag & Auf Basis tatsächlicher Kosten der Periode \\
Normal-Zuschlag & Durchschnitt mehrerer Perioden -- glättet Schwankungen \\
\hline
\end{tabular}

\end{multicols}

\newpage

% ============================================================
%  SEITE 5 -- KALKULATION (KOSTENTRÄGERRECHNUNG)
% ============================================================
\section{5 \quad Kostenträgerrechnung -- Kalkulation (3. Stufe der KLR)}

\begin{multicols}{2}

\subsection{Kalkulationsschema (Zuschlagskalkulation)}

\begin{formelbox}
\renewcommand{\arraystretch}{1.15}
\begin{tabular}{@{}lr@{}}
Materialeinzelkosten (MEK) & \\
$+$ Materialgemeinkosten: MEK $\times$ MGK\% & \\
\hline
$=$ \textbf{Materialkosten (MK)} & \\[2pt]
$+$ Fertigungslöhne (FL) & \\
$+$ Fertigungsgemeinkosten: FL $\times$ FGK\% & \\
$+$ Sondereinzelkosten Fertigung (SEF) & \\
\hline
$=$ \textbf{Fertigungskosten (FK)} & \\[2pt]
\hline
$=$ \textbf{Herstellkosten (HK)} $=$ MK $+$ FK & \\[2pt]
$+$ Verwaltungsgemeinkosten: HK $\times$ VwGK\% & \\
$+$ Vertriebsgemeinkosten: HK $\times$ VtGK\% & \\
$+$ Sondereinzelkosten Vertrieb (SEV) & \\
\hline
$=$ \textbf{Selbstkosten (SK)} & \\[2pt]
$+$ Gewinnzuschlag (\%) & \\
\hline
$=$ \textbf{Nettoverkaufspreis} & \\
$+$ MwSt. (19\%) & \\
\hline
$=$ \textbf{Bruttoverkaufspreis} & \\
\end{tabular}
\end{formelbox}

\subsection{Retrograde Kalkulation (Rückwärts)}

Vom Marktpreis rückwärts zu den maximal erlaubten MEK:

\begin{formelbox}
\renewcommand{\arraystretch}{1.1}
\begin{tabular}{@{}l@{}}
Bruttoverkaufspreis \\
$-$ MwSt. \\
$=$ Nettoangebotspreis \\
$-$ Gewinnzuschlag \\
$=$ Selbstkosten \\
$-$ VtGK, VwGK, SEV \\
$=$ Herstellkosten \\
$-$ FGK, FL, SEF \\
$=$ Materialkosten \\
$-$ MGK \\
$=$ \textbf{max. zulässige MEK} \\
\end{tabular}
\end{formelbox}

\columnbreak

\subsection{Divisionskalkulation}

Anwendung: \textbf{Einproduktunternehmen} / homogene Massenfertigung.

\begin{formelbox}
\[
k = \frac{K_{\text{gesamt}}}{x}
\]
\end{formelbox}

\begin{tabular}{|L{1.5cm}|L{5.5cm}|}
\hline $k$ & Stückkosten (€/Stück) \\
$K$ & Gesamtkosten der Periode \\
$x$ & Produktionsmenge \\
\hline
\end{tabular}

\vspace{4pt}
\textbf{Mehrstufige Divisionskalkulation} (mit Lagerbeständen):
\begin{formelbox}
\[
k = \frac{K_1}{x_1} + \frac{K_2}{x_2} + \ldots + \frac{K_n}{x_n}
\]
\end{formelbox}

\subsection{Äquivalenzziffernkalkulation}

Anwendung: \textbf{Mehrproduktunternehmen} mit ähnlichen Produkten (unterschiedlicher Ressourcenverbrauch).

\begin{formelbox}
\[
k_E = \frac{K_{\text{gesamt}}}{\sum_i (x_i \times \ddot{A}Z_i)}
\qquad
k_i = k_E \times \ddot{A}Z_i
\]
\end{formelbox}

\begin{tabular}{|L{2cm}|L{5cm}|}
\hline
\rowcolor{rowbg}
$k_E$ & Einheitsstückkosten (Kosten für Referenzprodukt mit $\ddot{A}Z=1$) \\
$\ddot{A}Z_i$ & Äquivalenzziffer des Produkts $i$ (relativer Kostenbedarf) \\
$x_i$ & Produktionsmenge des Produkts $i$ \\
$k_i$ & Stückkosten Produkt $i$ \\
\hline
\end{tabular}

\subsection{Prozesskostenkalkulation}

Anwendung: Gemeinkosten werden nach tatsächlicher Ressourcennutzung (Aktivitäten) verrechnet.

\begin{formelbox}
\[
\text{Prozesskostensatz} = \frac{\text{Prozesskosten}}{\text{Kostentreibermenge}}
\]
\[
\text{Prozesskosten/Stück} = \sum_j \frac{\text{Prozesskostensatz}_j \times n_j}{x}
\]
\end{formelbox}

\begin{tabular}{|L{2cm}|L{5cm}|}
\hline
$n_j$ & Anzahl der Aktivitäten des Prozesses $j$ pro Periode \\
$x$ & Produktionsmenge \\
\hline
\end{tabular}

\end{multicols}

\newpage

% ============================================================
%  SEITE 6 -- GEWINNVERTEILUNG & RECHTSFORMEN
% ============================================================
\section{6 \quad Gewinnverteilung und Rechtsformen}

\begin{multicols}{2}

\subsection{OHG -- Offene Handelsgesellschaft}

Alle Gesellschafter haften \textbf{unbeschränkt} persönlich.

\begin{formelbox}
\textbf{Schritt 1: Verzinsung der Kapitaleinlagen mit 4\%}\\[3pt]
$Z_i = K_i \times 4\%$\\[4pt]
\textbf{Schritt 2: Restgewinn gleichmäßig pro Kopf}\\[3pt]
$G_{\text{Rest}} = G - \sum Z_i$\\
$G_{\text{Rest,\,Kopf}} = G_{\text{Rest}} \div n$\\[4pt]
\textbf{Schritt 3: Gesamtanteil Gesellschafter $i$}\\[3pt]
$G_i = Z_i + G_{\text{Rest,\,Kopf}}$
\end{formelbox}

\begin{tabular}{|L{2cm}|L{5cm}|}
\hline
\rowcolor{rowbg}
$G$ & Jahresgewinn gesamt \\
$K_i$ & Kapitaleinlage von Gesellschafter $i$ \\
$n$ & Anzahl der Gesellschafter \\
$Z_i$ & Verzinsung für Gesellschafter $i$ \\
\hline
\end{tabular}

\subsection{KG -- Kommanditgesellschaft}

Komplementär (Kpl): unbeschränkt haftend; Kommanditist (Kdt): haftet nur mit Einlage.

\begin{formelbox}
\textbf{Schritt 1: Unternehmerlohn für Komplementär}\\[3pt]
\textbf{Schritt 2: Verzinsung aller Kapitaleinlagen}\\[3pt]
$Z_i = K_i \times i\%$\\[4pt]
\textbf{Schritt 3: Restgewinn nach Kapitalanteil}\\[3pt]
$G_{\text{Rest},i} = G_{\text{Rest}} \times \dfrac{K_i}{\sum_j K_j}$\\[4pt]
\textbf{Schritt 4: Gesamtanteil = Lohn + Zinsen + Restanteil}
\end{formelbox}

\columnbreak

\subsection{GmbH}

\begin{formelbox}
Gewinnverteilung nach Kapitalanteil (Stammeinlage):
\[
G_i = G_{\text{gesamt}} \times \frac{\text{Stammeinlage}_i}{\text{Stammkapital gesamt}}
\]
\end{formelbox}

\subsection{Rechtsformen im Überblick}

\begin{tabular}{|L{2.3cm}|L{2.5cm}|L{1.5cm}|L{0.8cm}|}
\hline
\rowcolor{rowbg}
\textbf{Rechtsform} & \textbf{Haftung} & \textbf{Mindest-kapital} & \textbf{GS} \\
\hline
Einzel\-unter\-nehmen & unbeschränkt persönl. & -- & 1 \\
\hline
GbR & unbeschränkt persönl. & -- & $\geq\!2$ \\
\hline
OHG & unbeschränkt persönl. & -- & $\geq\!2$ \\
\hline
KG & Kpl: unbeschr.\ Kdt: beschr. & -- & $\geq\!2$ \\
\hline
GmbH & nur Gesellschafts-vermögen & 25.000\,€ & $\geq\!1$ \\
\hline
AG & nur Gesellschafts-vermögen & 50.000\,€ & $\geq\!1$ \\
\hline
\end{tabular}

\subsection{Organisationsformen}

\begin{tabular}{|L{3cm}|L{4.2cm}|}
\hline
\rowcolor{rowbg}
\textbf{Organisationsform} & \textbf{Merkmal} \\
\hline
Einlinie & Ein Vorgesetzter je Mitarbeiter; klare Hierarchie \\
\hline
Mehrlinie & Mehrere Vorgesetzte (Fachspezialisten); Kompetenzstreitigkeiten möglich \\
\hline
Matrix & Fach- \& Objektverantwortung; Doppelunterstellung \\
\hline
Center-/Spartenorg. & Autonome Geschäftsbereiche mit eigenem Ergebnis (Profit Center) \\
\hline
\end{tabular}

\end{multicols}

\newpage

% ============================================================
%  SEITE 7 -- KOSTENFUNKTIONEN & VERFAHRENSVERGLEICH
% ============================================================
\section{7 \quad Kostenfunktionen, Verfahrensvergleich und Preisuntergrenze}

\begin{multicols}{2}

\subsection{Linearer Kostenverlauf}

\begin{formelbox}
\[
K(x) = K_f + k_v \cdot x \qquad \text{(Gesamtkosten)}
\]
\[
k(x) = \underbrace{\frac{K_f}{x}}_{\text{Stückfixkosten}} + k_v \qquad \text{(Stückkosten)}
\]
\[
E(x) = p \cdot x \qquad \text{(Erlösfunktion)}
\]
\end{formelbox}

\begin{tabular}{|L{1.8cm}|L{5.2cm}|}
\hline
\rowcolor{rowbg}
$K(x)$ & Gesamtkosten bei Menge $x$ \\
$K_f$ & Fixkosten (konstant, unabhängig von $x$) \\
$k_v$ & Variable Stückkosten (€/Stück, konstant) \\
$k(x)$ & Gesamtstückkosten (fallen mit steigendem $x$) \\
$E(x)$ & Gesamterlös (Umsatz) \\
$p$ & Verkaufspreis je Einheit \\
$x$ & Ausbringungsmenge \\
\hline
\end{tabular}

\vspace{4pt}
\textbf{$k_v$ und $K_f$ aus zwei Punkten ableiten:}
\begin{formelbox}
\[
k_v = \frac{K_2 - K_1}{x_2 - x_1}
\qquad
K_f = K_1 - k_v \cdot x_1
\]
\end{formelbox}

\subsection{Verfahrensvergleich (kritische Menge)}

Zwei Fertigungsverfahren mit unterschiedlichen Kosten:

\begin{formelbox}
\[
x_{\text{krit}} = \frac{K_{f2} - K_{f1}}{k_{v1} - k_{v2}}
\qquad (\text{mit } K_{f2} > K_{f1},\; k_{v1} > k_{v2})
\]
\end{formelbox}

\begin{itemize}[noitemsep, leftmargin=*]
  \item $x < x_{\text{krit}}$: Verfahren mit \textbf{niedrigeren Fixkosten} (Verfahren 1) günstiger
  \item $x > x_{\text{krit}}$: Verfahren mit \textbf{niedrigeren variablen Kosten} (Verfahren 2) günstiger
\end{itemize}

\columnbreak

\subsection{Break-Even-Point (Gewinnschwelle)}

\begin{formelbox}
\[
E = K \;\Rightarrow\; p \cdot x_{\text{BEP}} = K_f + k_v \cdot x_{\text{BEP}}
\]
\[
\boxed{x_{\text{BEP}} = \frac{K_f}{p - k_v} = \frac{K_f}{db}}
\]
\end{formelbox}

$db = p - k_v$ = Deckungsbeitrag je Stück\\
Bei $x = x_{\text{BEP}}$: Gewinn = 0. Erst darüber wird Gewinn erzielt.

\subsection{Preisuntergrenze}

\begin{tabular}{|L{2.5cm}|L{5cm}|}
\hline
\rowcolor{rowbg}
\textbf{Fristigkeit} & \textbf{Preisuntergrenze} \\
\hline
\textbf{Kurzfristig} & $p_{\min} = k_v$ (var. Stückkosten)\\
& Bedingung: $db = p - k_v \geq 0$\\
& (Fixkosten sind bereits angefallen) \\
\hline
\textbf{Langfristig} & $p_{\min} = k$ (Selbstkosten/Stück)\\
& Alle Kosten inkl. Fixkosten müssen gedeckt sein \\
\hline
\end{tabular}

\subsection{Nichtlineare Kostenverläufe}

\begin{tabular}{|L{2.8cm}|L{4.5cm}|}
\hline
\rowcolor{rowbg}
\textbf{Verlauf} & \textbf{Merkmal} \\
\hline
Degressiv (unterproportional) & Grenzkosten sinken mit steigendem $x$ (z.\,B. Mengenrabatte, Lernkurve) \\
Progressiv (überproportional) & Grenzkosten steigen (z.\,B. Überstundenzuschläge, Maschinenüberlastung) \\
Ertragsgesetzlich & Erst degressiv, dann progressiv; Wendepunkt = minimale Grenzkosten \\
\hline
\end{tabular}

\vspace{4pt}
\begin{hinweisbox}
\textbf{Kostenremanenz:} Bei sinkender Beschäftigung bauen sich Fixkosten nicht sofort ab (Kündigungsfristen, langfristige Verträge). Die Kosten reagieren träge nach unten -- Kostensenkung hinkt der Beschäftigungsabnahme nach.
\end{hinweisbox}

\end{multicols}

\newpage

% ============================================================
%  SEITE 8 -- DECKUNGSBEITRAGSRECHNUNG
% ============================================================
\section{8 \quad Deckungsbeitragsrechnung (Teilkostenrechnung)}

\begin{multicols}{2}

\subsection{Einstufige Deckungsbeitragsrechnung}

\begin{formelbox}
\renewcommand{\arraystretch}{1.15}
\begin{tabular}{@{}lr@{}}
Erlös & $E = p \cdot x$ \\
$-$ Variable Kosten & $K_v = k_v \cdot x$ \\
\hline
$=$ \textbf{Deckungsbeitrag} & $\mathbf{DB = E - K_v}$ \\
$-$ Fixkosten & $K_f$ \\
\hline
$=$ \textbf{Gewinn/Verlust} & $G = DB - K_f$ \\
\end{tabular}
\end{formelbox}

\textbf{Wichtige Formeln:}
\begin{formelbox}
\[
db = p - k_v \qquad \text{(Stückdeckungsbeitrag)}
\]
\[
DB = db \cdot x \qquad \text{(Gesamt-DB)}
\]
\[
DBR = \frac{db}{p} = \frac{DB}{E} \qquad \text{(Deckungsbeitragsrate)}
\]
\[
x_{\text{BEP}} = \frac{K_f}{db} \qquad \text{(Break-Even-Menge)}
\]
\end{formelbox}

\begin{tabular}{|L{2cm}|L{5cm}|}
\hline
\rowcolor{rowbg}
$db$ & Deckungsbeitrag je Stück (€/Stk.) -- Beitrag zur Fixkostendeckung \\
$DB$ & Gesamt-Deckungsbeitrag (€) \\
$DBR$ & Deckungsbeitragsrate -- Anteil des DB am Erlös \\
$K_v$ & Gesamte variable Kosten \\
$K_f$ & Gesamte Fixkosten (periodengebunden) \\
\hline
\end{tabular}

\columnbreak

\subsection{Mehrstufige Deckungsbeitragsrechnung}

Fixkosten werden in Schichten aufgeteilt (Produkt-, Gruppen-, Bereichs-, Unternehmensebene). Zeigt, welche Ebene tatsächlich rentabel ist.

\begin{formelbox}
\renewcommand{\arraystretch}{1.15}
\begin{tabular}{@{}l@{}}
Erlöse aller Produkte \\
$-$ variable Kosten \\
\hline
$=$ DB\,I (Stückbeitrag) \\
$-$ Produktfixkosten (produktspezifisch) \\
\hline
$=$ DB\,II (Produktgruppenbeitrag) \\
$-$ Produktgruppen-Fixkosten \\
\hline
$=$ DB\,III (Bereichsbeitrag) \\
$-$ Bereichsfixkosten \\
\hline
$=$ DB\,IV (Unternehmensbeitrag) \\
$-$ Unternehmensfixkosten \\
\hline
$=$ \textbf{Ergebnis (Gewinn / Verlust)} \\
\end{tabular}
\end{formelbox}

\textbf{Vorteil:} Produkt mit negativem DB\,I ist schädlich; mit negativem DB\,II aber noch rentabel, wenn Fixkosten bei Streichung nicht wegfallen.

\subsection{Relativer Deckungsbeitrag (Engpassrechnung)}

Bei einem \textbf{Engpass} (knappe Ressource, z.\,B. Maschinenzeit) Produktionsprogramm nach maximalem DB je Engpasseinheit optimieren.

\begin{formelbox}
\[
db_{\text{rel}} = \frac{db}{t} \qquad \left[\frac{\text{€}}{\text{Minute}}\right]
\]
\end{formelbox}

$t$ = Engpasszeit je Stück auf der Engpassmaschine (in Minuten)

\vspace{4pt}
\textbf{Vorgehensweise:}
\begin{enumerate}[noitemsep, leftmargin=*]
  \item $db_{\text{rel}}$ für jedes Produkt berechnen
  \item Produkte absteigend nach $db_{\text{rel}}$ sortieren
  \item Engpasskapazität zuerst dem Produkt mit höchstem $db_{\text{rel}}$ zuteilen
  \item Weiter bis Kapazität erschöpft; letztes Produkt evtl. anteilig
\end{enumerate}

\begin{hinweisbox}
\textbf{Wichtig:} Bei \textbf{keinem} Engpass gilt: Alle Produkte produzieren, die $db > 0$ haben. Erst bei Engpass zählt $db_{\text{rel}}$.
\end{hinweisbox}

\end{multicols}

\newpage

% ============================================================
%  SEITE 9 -- BILANZ & BILANZANALYSE
% ============================================================
\section{9 \quad Bilanz und Bilanzanalyse}

\begin{multicols}{2}

\subsection{Bilanzstruktur (vereinfacht)}

\begin{tabular}{|L{3.5cm}|L{3.5cm}|}
\hline
\rowcolor{rowbg}
\multicolumn{1}{|c|}{\textbf{AKTIVA}} & \multicolumn{1}{c|}{\textbf{PASSIVA}} \\
\hline
\textbf{Anlagevermögen (AV)} & \textbf{Eigenkapital (EK)} \\
\quad Sachanlagen & \quad Gezeichnetes Kapital \\
\quad Immaterielle AV & \quad Rücklagen \\
\quad Finanzanlagen & \quad Jahresüberschuss \\
\hline
\textbf{Umlaufvermögen (UV)} & \textbf{Fremdkapital (FK)} \\
\quad Vorräte & \quad Rückstellungen \\
\quad Forderungen & \quad Langfristige Verb. ($>$1\,J.) \\
\quad Wertpapiere & \quad Kurzfristige Verb. ($<$1\,J.) \\
\quad Flüssige Mittel & \\
\hline
Aktive RAP & Passive RAP \\
\hline
\textbf{Bilanzsumme} = & \textbf{Bilanzsumme} \\
\multicolumn{2}{|c|}{$\sum \text{Aktiva} = \sum \text{Passiva}$ \quad (immer!)} \\
\hline
\end{tabular}

\subsection{EBIT-Berechnung}

\begin{formelbox}
\renewcommand{\arraystretch}{1.1}
\begin{tabular}{@{}l@{}}
Jahresüberschuss / -fehlbetrag \\
$\pm$ Außerordentliches Ergebnis \\
$+$ Steueraufwand (Ertragsteuern) \\
$\pm$ Finanzergebnis (Zinsaufwand $-$ Zinsertrag) \\
\hline
$=$ \textbf{EBIT} \\
\quad (Earnings Before Interest and Taxes) \\
\end{tabular}
\end{formelbox}

EBIT = Betriebsergebnis, unabhängig von Finanzierung und Steuern.

\columnbreak

\subsection{Kennzahlen der Bilanzanalyse}

\textbf{Liquiditätskennzahlen:}

\begin{formelbox}
\[
L_1 = \frac{\text{Flüssige Mittel}}{\text{kurzfristige Verbindlichkeiten}} \times 100\%
\]
\[
L_2 = \frac{\text{Flüssige Mittel} + \text{kurzfristige Forderungen}}{\text{kurzfristige Verbindlichkeiten}} \times 100\%
\]
\[
L_3 = \frac{\text{Umlaufvermögen}}{\text{kurzfristige Verbindlichkeiten}} \times 100\%
\]
\end{formelbox}

\textbf{Richtwerte:} $L_1 \approx 20\%$; $L_2 \approx 100\%$; $L_3 \approx 200\%$

\vspace{4pt}
\textbf{Kapital- und Ertragskennzahlen:}

\begin{formelbox}
\[
\text{Verschuldungsgrad} = \frac{FK}{EK} \times 100\%
\]
\[
\text{Eigenkapitalquote} = \frac{EK}{\text{Gesamtkapital}} \times 100\%
\]
\[
\text{GK-Rendite} = \frac{\text{Gewinn} + \text{FK-Zinsen}}{\text{Gesamtkapital}} \times 100\%
\]
\[
\text{EK-Rendite} = \frac{\text{Gewinn}}{EK} \times 100\%
\]
\[
\text{Umsatzrendite} = \frac{\text{Gewinn}}{\text{Umsatz}} \times 100\%
\]
\end{formelbox}

\subsection{Du-Pont-Kennzahlensystem (ROI)}

\begin{formelbox}
\[
ROI = \underbrace{\frac{\text{Gewinn}}{\text{Umsatz}}}_{\text{Umsatzrendite}} \times \underbrace{\frac{\text{Umsatz}}{\text{Gesamtkapital}}}_{\text{Kapitalumschlag}}
= \frac{\text{Gewinn}}{\text{Gesamtkapital}}
\]
\end{formelbox}

\textbf{Stellhebel zur ROI-Steigerung:}
\begin{itemize}[noitemsep, leftmargin=*]
  \item Umsatzrendite: Kosten senken oder Preise erhöhen
  \item Kapitalumschlag: Kapital effizienter einsetzen, Lagerbestände reduzieren
\end{itemize}

\end{multicols}

\newpage

% ============================================================
%  SEITE 10 -- VWL GRUNDBEGRIFFE & KREISLAUF/BIP
% ============================================================
\section{10 \quad VWL: Grundbegriffe und Volkseinkommen (BIP)}

\begin{multicols}{2}

\subsection{Wirtschaftliche Grundbegriffe}

\begin{vwlbox}
\[
\text{Wirtschaftlichkeit} = \frac{\text{Leistung (Output)}}{\text{Kosten (Input)}} \;\geq 1
\]
\[
\text{Produktivität} = \frac{\text{Produktionsmenge (Output)}}{\text{Faktoreinsatz (Input)}}
\]
\[
\text{Umsatzrendite} = \frac{\text{Gewinn}}{\text{Umsatz}} \times 100\%
\]
\[
\text{Kapitalrendite} = \frac{\text{Gewinn}}{\text{eingesetztes Kapital}} \times 100\%
\]
\end{vwlbox}

\textbf{Ökonomisches Prinzip:}

\begin{tabular}{|L{2.5cm}|L{4.5cm}|}
\hline
\rowcolor{rowvwl}
\textbf{Prinzip} & \textbf{Ziel} \\
\hline
Maximalprinzip & Mit gegebenem Input $\to$ maximalen Output erreichen \\
\hline
Minimalprinzip & Bestimmten Output $\to$ mit minimalem Input erreichen \\
\hline
\end{tabular}

\subsection{Wirtschaftskreislauf}

\textbf{Einfacher Kreislauf} (2 Sektoren):
\begin{itemize}[noitemsep, leftmargin=*]
  \item Haushalte $\leftrightarrow$ Unternehmen
  \item Faktoren (Arbeit/Kapital/Boden) gegen Einkommen; Güter gegen Konsum
\end{itemize}

\textbf{Erweiterter Kreislauf} (+\,Banken/Staat):
\begin{itemize}[noitemsep, leftmargin=*]
  \item Banken: Vermittlung Spareinlagen zu Investitionskrediten
  \item Staat: Steuern, Transferzahlungen, öffentliche Ausgaben
\end{itemize}

\textbf{Vollständiger Kreislauf} (5 Sektoren: Haushalte, Unternehmen, Banken, Staat, Ausland)

\columnbreak

\subsection{BIP -- Berechnungsmethoden}

\textbf{Entstehungsrechnung:}
\begin{vwlbox}
\[
\text{Produktionswert} - \text{Vorleistungen} = \text{Bruttowertschöpfung}
\]
\[
+ \text{Gütersteuern} - \text{Gütersubventionen} = \mathbf{BIP}
\]
\end{vwlbox}

\textbf{Verwendungsrechnung:}
\begin{vwlbox}
\renewcommand{\arraystretch}{1.1}
\begin{tabular}{@{}l@{}}
Privater Konsum ($C$) \\
$+$ Staatlicher Konsum ($G$) \\
$+$ Bruttoinvestitionen ($I$) \\
$+$ Außenbeitrag (Exporte $-$ Importe) \\
\hline
$=$ \textbf{BIP} \\
\end{tabular}
\end{vwlbox}

\textbf{Verteilungsrechnung:}
\begin{vwlbox}
\renewcommand{\arraystretch}{1.1}
\begin{tabular}{@{}l@{}}
Arbeitnehmerentgelt \\
$+$ Unternehmens- und Vermögenseinkommen \\
\hline
$=$ Volkseinkommen (NSOz.Fakt.) \\
$+$ Produktions- und Importabgaben \\
$-$ Subventionen + Abschreibungen + Saldo \\
\hline
$=$ \textbf{BIP} \\
\end{tabular}
\end{vwlbox}

\textbf{Weitere BIP-Formeln:}
\begin{vwlbox}
\[
\text{BIP}_{\text{real}} = \frac{\text{BIP}_{\text{nominal}}}{\text{Preisindex}} \times 100
\]
\begin{align*}
\text{BNE} &= \text{BIP} - \text{Inlandseinkommen Ausländer}\\
           &\quad + \text{Auslandseinkommen Inländer}
\end{align*}
\begin{align*}
\text{NEW} &= \text{BIP} - \text{soziale Kosten}\\
           &\quad + \text{private Dienste (unentgeltlich)}
\end{align*}
\end{vwlbox}

\begin{tabular}{|L{2cm}|L{5cm}|}
\hline
\rowcolor{rowvwl}
BNE & Bruttonationaleinkommen (nach Inländerprinzip) \\
NEW & Net Economic Welfare (Nettowohlstandsindikator) \\
\hline
\end{tabular}

\end{multicols}

\newpage

% ============================================================
%  SEITE 11 -- GELD, GELDMENGE & GELDPOLITIK
% ============================================================
\section{11 \quad Geldmengen, Geldschöpfung und EZB-Instrumente}

\begin{multicols}{2}

\subsection{Geldmengenaggregate der EZB}

\begin{tabular}{|L{1cm}|L{6cm}|}
\hline
\rowcolor{rowvwl}
\textbf{M0} & Zentralbankgeld: Bargeld im Umlauf + Mindestreserveguthaben der Banken bei EZB (Geldbasis) \\
\hline
\textbf{M1} & M0 $+$ Sichteinlagen (täglich fällige Girokonten) \\
\hline
\textbf{M2} & M1 $+$ Einlagen mit Laufzeit $\leq 2$ Jahre oder Kündigungsfrist $\leq 3$ Monate \\
\hline
\textbf{M3} & M2 $+$ Geldmarktfondsanteile $+$ Repogeschäfte $+$ Bankschuldverschreibungen ($\leq 2$ J.) \\
\hline
\end{tabular}

\textbf{Merkhilfe:} M1 $\subset$ M2 $\subset$ M3 (jede Menge enthält die vorherige vollständig)

\subsection{Geldschöpfungsmultiplikator}

Geschäftsbanken können durch Kreditvergabe neues Buchgeld schaffen (Giralgeldschöpfung).

\begin{vwlbox}
\[
m = \frac{1}{r + b}
\]
\[
\Delta M = m \times \Delta Z\!B\text{-}{\text{Geld}}
\]
\end{vwlbox}

\begin{tabular}{|L{1.5cm}|L{5.5cm}|}
\hline
\rowcolor{rowvwl}
$m$ & Geldschöpfungsmultiplikator \\
$r$ & Mindestreservesatz (MRS): Pflichtanteil der Einlagen bei EZB \\
$b$ & Barreservesatz (BRS): Anteil, den Banken als Bargeld im Tresor halten \\
$\Delta M$ & Gesamte Geldmengenausweitung durch Multiplikation \\
\hline
\end{tabular}

\columnbreak

\subsection{Tenderverfahren der EZB}

\textbf{Mengentender} (Festzinssatz):
\begin{vwlbox}
\[
\text{Repartierungsquote} = \frac{\text{Zuteilungsbetrag}}{\sum \text{Gebote aller Banken}} \times 100\%
\]
\end{vwlbox}
EZB legt Zinssatz und Zuteilungsbetrag fest. Alle Banken erhalten denselben prozentualen Anteil ihrer Gebote.

\vspace{4pt}
\textbf{Zinstender} (variabler Zinssatz):
\begin{itemize}[noitemsep, leftmargin=*]
  \item Banken bieten Zinssatz und Betrag
  \item Höchste Zinssätze werden zuerst bedient
  \item EZB legt Mindestbietungssatz $=$ Leitzins fest
\end{itemize}

\subsection{EZB-Leitzinsen und Geldpolitikinstrumente}

\begin{tabular}{|L{3.5cm}|L{4cm}|}
\hline
\rowcolor{rowvwl}
\textbf{Instrument} & \textbf{Funktion} \\
\hline
Haupt\-refi\-nan\-zie\-rungs\-geschäft & Wochentender; Standardinstrument; Leitzins \\
\hline
Spitzen\-refi\-nan\-zie\-rungs\-fazilität & Übernachtkredite; obere Zinsgrenze (Ceiling) \\
\hline
Einlagefazilität & Übernachteinlagen bei EZB; untere Zinsgrenze (Floor) \\
\hline
Offenmarktgeschäfte & Kauf/Verkauf von Wertpapieren am offenen Markt (kurz- bis langfristig) \\
\hline
Mindestreservepolitik & Pflichtreservesatz senken $\to$ mehr Kreditvergabe möglich \\
\hline
\end{tabular}

\vspace{4pt}
\textbf{Wirkungskette Leitzinserhöhung (kontraktive Geldpolitik):}
\begin{itemize}[noitemsep, leftmargin=*]
  \item $\uparrow$ Leitzins $\to$ $\uparrow$ Refinanzierungskosten der Banken
  \item $\to$ $\uparrow$ Kreditzinsen $\to$ $\downarrow$ Kreditnachfrage
  \item $\to$ $\downarrow$ Investitionen, $\downarrow$ Konsum $\to$ $\downarrow$ Inflation
\end{itemize}

\textbf{Wirkungskette Leitzinssenkung (expansive Geldpolitik):}
\begin{itemize}[noitemsep, leftmargin=*]
  \item $\downarrow$ Leitzins $\to$ $\downarrow$ Kreditkosten $\to$ $\uparrow$ Kredite $\to$ $\uparrow$ Investitionen $\to$ $\uparrow$ BIP
\end{itemize}

\end{multicols}

\newpage

% ============================================================
%  SEITE 12 -- INFLATION, MARKTFORMEN & PREISELASTIZITÄT
% ============================================================
\section{12 \quad Inflation, Marktformen und Preiselastizität}

\begin{multicols}{2}

\subsection{Inflation}

\begin{vwlbox}
\[
\text{Inflationsrate} = \frac{P_t - P_{t-1}}{P_{t-1}} \times 100\%
\]
\end{vwlbox}

\begin{tabular}{|L{3.5cm}|L{3.5cm}|}
\hline
\rowcolor{rowvwl}
\textbf{Art} & \textbf{Rate / Merkmal} \\
\hline
Schleichende Inflation & $< 5\%$ p.\,a. -- kaum spürbar \\
Trabende Inflation & $5\text{--}20\%$ p.\,a. \\
Galoppierende Inflation & $> 20\%$ p.\,a. -- wirtschaftsschädigend \\
Hyperinflation & $> 50\%$ pro Monat (nach Cagan) \\
\hline
Offene Inflation & Preise steigen tatsächlich \\
Verdeckte Inflation & Preise staatlich gedeckelt; Güterknappheit und Schwarzmarkt \\
\hline
\end{tabular}

\vspace{4pt}
\textbf{Binnenwert des Geldes:} Kaufkraft im Inland -- sinkt bei Inflation \\
\textbf{Außenwert des Geldes:} Wechselkurs -- Verhältnis zu Fremdwährungen

\subsection{Marktformmatrix}

\begin{tabular}{|C{1.8cm}|C{1.5cm}|C{1.5cm}|C{1.5cm}|}
\hline
\rowcolor{rowvwl}
& \multicolumn{3}{c|}{\textbf{Anbieter}} \\
\rowcolor{rowvwl}
\textbf{Nachfrager} & Viele & Wenige & Einer \\
\hline
\textbf{Viele} & Polypol & Oligopol & Monopol \\
\textbf{Wenige} & Nachfrage-oligopol & Bilat. Oligopol & Beschr. Monopol \\
\textbf{Einer} & Monopson & Beschr. Monopson & Bilat. Monopol \\
\hline
\end{tabular}

\columnbreak

\subsection{Preiselastizität der Nachfrage}

\begin{vwlbox}
\[
\varepsilon = \frac{\Delta Q / Q}{\Delta P / P}
= \frac{\Delta Q}{\Delta P} \times \frac{P}{Q}
\]
\end{vwlbox}

\begin{tabular}{|L{1.8cm}|L{2.2cm}|L{3cm}|}
\hline
\rowcolor{rowvwl}
\textbf{Wert $|\varepsilon|$} & \textbf{Bezeichnung} & \textbf{Bedeutung} \\
\hline
$\to \infty$ & Vollkommen elastisch & Kleinste Preisänderung $\Rightarrow$ unendliche Mengenänderung (z.\,B. homog. Markt) \\
\hline
$> 1$ & Elastisch & \%\,Mengenänd. $>$ \%\,Preisänd. (Luxusgüter, Substitute vorhanden) \\
\hline
$= 1$ & Proportional (unitär) & \%\,Mengenänd. $=$ \%\,Preisänd. \\
\hline
$< 1$ & Unelastisch & \%\,Mengenänd. $<$ \%\,Preisänd. (Grundbedarf, keine Substitute) \\
\hline
$= 0$ & Vollkommen unelastisch & Preisänderung hat keinen Einfluss auf Menge \\
\hline
$> 0$ (positiv) & Anormal (Giffen) & Preiserhöhung $\Rightarrow$ mehr Nachfrage (inferiore Güter) \\
\hline
\end{tabular}

\vspace{4pt}
\begin{tabular}{|L{2cm}|L{5cm}|}
\hline
\rowcolor{rowvwl}
$\varepsilon$ & Preiselastizitätskoeffizient \\
$\Delta Q$ & Absolute Änderung der Nachfragemenge \\
$Q$ & Ausgangsmenge \\
$\Delta P$ & Absolute Preisänderung \\
$P$ & Ausgangspreis \\
\hline
\end{tabular}

\vspace{4pt}
\textbf{Umsatzwirkung einer Preiserhöhung:}\par
\begin{tabular}{|L{3cm}|L{4cm}|}
\hline
$|\varepsilon| > 1$ (elastisch) & Preis $\uparrow$ $\Rightarrow$ Umsatz $\downarrow$ \\
$|\varepsilon| < 1$ (unelastisch) & Preis $\uparrow$ $\Rightarrow$ Umsatz $\uparrow$ \\
$|\varepsilon| = 1$ (unitär) & Preis $\uparrow$ $\Rightarrow$ konstant \\
\hline
\end{tabular}

\begin{hinweisbox}
\textbf{Vorzeichen:} Preiselastizität ist normalerweise \textbf{negativ} (Preis $\uparrow$ $\to$ Menge $\downarrow$). Bei Giffen-Gütern positiv. In der Praxis oft als Absolutbetrag $|\varepsilon|$ angegeben.
\end{hinweisbox}

\end{multicols}

\vspace{6pt}
{\color{bwlblue}\hrule height 1pt}
\vspace{3pt}
{\small\centering\textbf{Ende der Formelsammlung} -- Alle Formeln sind prüfungsrelevant. Schwerpunkt Klausur: 3 Stufen der KLR (Seiten 1--5), Deckungsbeitragsrechnung (Seite 8), BIP (Seite 10), Geldpolitik (Seite 11).}

\end{document}
